\section{Introducción}
\label{sec:introduccion-ch}

La hipótesis del continuo pregunta si existe un conjunto de números reales
cuya cardinalidad sea estrictamente mayor que la de los naturales y menor que
la de la recta completa. Este trabajo obtiene una respuesta negativa a esa
existencia, y por tanto afirmativa a la hipótesis, en una clausura genealógica
definida mediante las operaciones de Holografía Modular Triádica (HMT).
Su contribución enlaza la construcción de la recta por prolongaciones
compatibles con la comparación cardinal de todos sus subconjuntos internos.

El núcleo formal HMT reúne tres objetos que se desarrollan desde sus
definiciones. All Possible Paths (APP) dota al soporte
\((\mathbb Z/9\mathbb Z)^2\) de las evaluaciones de suma y producto, con sus
residuos y cocientes. Las traslaciones por \((\pm1,0)\) y \((0,\pm1)\)
determinan su grafo de Cayley; las celdas cuadradas y la identificación de
bordes opuestos le dan una realización toroidal bidimensional. TRIT designa
la firma ternaria orientada \(\{-1,0,+1\}\), con la que se distinguen los
regímenes de las transiciones. El Topological Phase Kernel (TPK), o núcleo
topológico de fase, compone selección, transporte y memoria para prolongar
los estados. Estas funciones son distintas y actúan conjuntamente: el soporte
sitúa los caminos, las evaluaciones publican su contenido aritmético y la
dinámica conserva la historia de cada transición.

El argumento sigue esa composición desde las operaciones finitas hasta el
dominio en el que se efectúa la comparación cardinal:
\[
 \mathrm{APP}\longrightarrow\mathrm{TRIT}\longrightarrow\mathrm{TPK}
 \longrightarrow X_\infty^{\mathrm{enr}}
 \longrightarrow\mathcal C_{\mathrm{cont}}^{\mathrm{disc}}
 \longrightarrow m_\infty
 \longrightarrow\mathfrak G_{\mathrm{HMT}}(h,m_\infty).
\]
La composición produce estados enriquecidos, prolongaciones compatibles,
historias completas y, sólo después, una realización arquimediana. La recta
real es salida de la construcción, no su universo de partida; la cuestión
cardinal se decide sobre el dominio íntegro producido por esa genealogía.

La cuestión central es si la recta publicada por el sistema admite, dentro de
su clausura genealógica completa, un cardinal estrictamente intermedio entre
\(\aleph_0\) y el cardinal de la propia recta. El teorema se formula en
\(\mathfrak G_{\mathrm{HMT}}(h,m_\infty)\), la clausura generada por el sistema
completo. En ese dominio cuantifica sobre
\(\mathcal P^{\mathfrak G}(\mathbb R_{\mathrm{HMT}}^{\mathfrak G})\) y
establece la dicotomía cardinal completa.

El teorema~\ref{thm:decision-continuo-hmt}, demostrado en la
sección~\ref{sec:cierre-cardinal-nativo}, establece el resultado principal.
En las comparaciones cardinales, la recta y los conjuntos potencia se
interpretan en la clausura genealógica
\(\mathfrak G=\mathfrak G_{\mathrm{HMT}}(h,m_\infty)\). Este dominio y sus
cuantificadores se construyen explícitamente antes de efectuar la comparación.
El resultado es el siguiente. Para toda historia realizada \(m_\infty\)
producida por el sistema completo asociado a
\[
 \mathrm{APP}\longrightarrow\mathrm{TRIT}\longrightarrow\mathrm{TPK}
 \longrightarrow X_\infty^{\mathrm{enr}}
 \longrightarrow\mathcal C_{\mathrm{cont}}^{\mathrm{disc}},
\]
cuyas restricciones son compatibles y cuyas prolongaciones son no vacías,
exhaustivas y supervivientes a todo horizonte, la extensión genealógica
\(\mathfrak G_{\mathrm{HMT}}(h,m_\infty)\) genera la recta integral,
cuantifica sobre toda su potencia y satisface
\[
 \mathfrak G_{\mathrm{HMT}}(h,m_\infty)
 \models 2^{\aleph_0}=\aleph_1.
\]
En particular, para toda
\(A\in\mathcal P^{\mathfrak G}(\mathbb R_{\mathrm{HMT}}^{\mathfrak G})\),
\[
 |A|\leq\aleph_0
 \qquad\text{o}\qquad
 |A|=|\mathbb R_{\mathrm{HMT}}^{\mathfrak G}|.
\]
Por tanto, HMT demuestra afirmativamente la hipótesis del continuo para el
continuo íntegro que define y genera. El enunciado formal y su prueba se dan
una sola vez en el teorema principal.

La diferencia esencial respecto de una presentación puramente extensional
reside en la identidad del objeto matemático. Una coordenada real aislada
conserva un valor; un estado HMT conserva además las operaciones que lo
produjeron, la hoja aritmética, la orientación trítica, los residuos y
cocientes atravesados, el acarreo, la frontera, la ruta y la memoria acumulada.
La proyección arquimediana no define ni agota esa estructura: publica una de
sus coordenadas. La segunda proyección conserva la incidencia de frontera que,
en desarrollos descendentes, alcanza Paley--Witt--Golay--Leech,
\(V^\natural\) y el Monstruo. Ambas realizaciones proceden del mismo estado y
ninguna se introduce para seleccionar retrospectivamente la otra.

\subsection{Construcción del continuo mediante prolongaciones compatibles}

APP dispone sobre un mismo soporte las hojas de suma y producto y conserva la
diferencia entre sus publicaciones. TRIT incorpora orientación y régimen. TPK
selecciona, transporta, memoriza y prolonga las transiciones admisibles. La
composición no termina en un estado finito: define un sistema de restricciones
\[
 \rho_{N+1,N}:X_{N+1}\longrightarrow X_N
\]
y extensiones no vacías
\[
 \operatorname{Ext}_N(x)\subseteq\rho_{N+1,N}^{-1}(x).
\]
La compatibilidad a toda profundidad produce el límite inverso
\[
 X_\infty^{\mathrm{enr}}
 =\varprojlim_N(X_N,\rho_{N+1,N}).
\]
El cuantificador es \(\forall N\); ninguna ejecución a mil, diez mil o un
millón de cifras constituye la cota del teorema.

Cada bloque sucesor subdivide un cilindro en \(729=3^6\) subcilindros de refinamiento disjuntos,
exhaustivos y de diámetro decreciente. Esta cifra describe la ramificación
local del refinamiento; no se identifica con el cardinal del continuo. La
cinta completa de bloques determina una historia y su proyección aritmética
determina un único punto real. Recíprocamente, la expansión canónica de cada
real reconstruye una cinta y la propiedad de cobertura la levanta a una
historia HMT compatible. Así, la recta es publicación sobreyectiva del árbol
construido, nunca un conjunto exterior empleado para escoger sus ramas.

La medición interviene como acoplamiento. Sobre el dominio considerado se
impone primero la condición de separación de fibras: el lector enriquecido
debe distinguir dos estados diferentes que publican el mismo valor. Bajo esa
condición, formalizada en la
proposición~\ref{prop:ch-recuperabilidad-fibra}, el instrumento distingue una
continuación compatible, añade al registro genealógico el resultado y la
retroacción y restringe las historias futuras a las que prolongan el nuevo
estado. La historia realizada \(m_\infty\) se forma mediante esas
actualizaciones. La transformación inversa reconstruye entonces el interior
desde la memoria trasladada a la frontera, de modo que la actualización no
equivale a pérdida de información. La
proposición~\ref{prop:ch-naturalidad-medida} muestra cuándo los cambios de
representación conservan también la medida de historias.

\subsection{Clausura semántica genealógica}

Sea \(h\) el código efectivo de los grafos de APP, TRIT, TPK, las
prolongaciones, las restricciones y las dos proyecciones. El par
\((h,m_\infty)\) se codifica primitivamente mediante un único real
\(u_{h,m}=h\oplus m_\infty\). Se define
\[
 \begin{aligned}
 \mathfrak G_0(h,m_\infty)
   &:=\operatorname{TC}\{\omega,u_{h,m},h,m_\infty\},\\
 \mathfrak G_{\beta+1}(h,m_\infty)
   &:=\operatorname{Def}_{\Sigma_{\mathrm{HMT}}}
        \bigl(\mathfrak G_\beta(h,m_\infty)\bigr),\\
 \mathfrak G_\lambda(h,m_\infty)
   &:=\bigcup_{\beta<\lambda}\mathfrak G_\beta(h,m_\infty),\\
 \mathfrak G_{\mathrm{HMT}}(h,m_\infty)
   &:=\bigcup_{\beta\in\operatorname{Ord}}
        \mathfrak G_\beta(h,m_\infty).
 \end{aligned}
\]
Esta clausura se formula mediante definibilidad relativa y condensación en el
sentido preciso de las jerarquías constructibles~\cite{Jech2003}. La hipótesis
del continuo no se introduce entre las condiciones que definen la clausura; se
obtiene mediante la demostración desarrollada en el artículo. Tampoco se
introducen como datos una enumeración de la recta ni una biyección objetivo. Su
operación sucesora publica los subconjuntos definibles mediante fórmulas finitas
del lenguaje HMT y parámetros ya producidos.
Equivalentemente, todo objeto admitido posee un árbol de derivación y todo
árbol admisible publica un objeto.

La jerarquía resultante es transitiva, contiene sus ordinales, satisface los
axiomas empleados en la comparación cardinal y posee un buen orden global
determinado por el primer nivel de nacimiento y el menor nombre de derivación.
La notación
\[
 \mathfrak G_{\mathrm{HMT}}(h,m_\infty)
 =L[u_{h,m}]=L[h,m_\infty]
\]
se obtiene después como teorema de representación. No constituye el motor de
APP, TRIT o TPK, no selecciona \(m_\infty\) y no sustituye la construcción
material del continuo.

\subsection{Comparación cardinal entre la recta HMT y el primer ordinal no numerable de la clausura}

La prueba cardinal se divide en dos direcciones independientes de la igualdad
buscada. Primero, la condensación genealógica demuestra que cada real generado
aparece antes de \(\omega_1^{\mathfrak G}\). Para mantener la misma
convención que en la prueba, sea \(\iota(r)\) el código del corte racional
de \(r\), sea \(\beta(r)\) el menor ordinal tal que
\(\iota(r)\in\mathfrak G_{\beta(r)+1}\), y sea \(n(r)\) el menor índice
que publica ese corte en la enumeración \(e_{\beta(r)+1}\). Se define
\[
 j_{\mathrm{HMT}}(r)=\omega\beta(r)+n(r),
 \qquad
 j_{\mathrm{HMT}}:
 \mathbb R_{\mathrm{HMT}}^{\mathfrak G}
 \hookrightarrow\omega_1^{\mathfrak G}.
\]

En la dirección opuesta, cada ordinal
\(\beta<\omega_1^{\mathfrak G}\) posee un código real \(w_\beta\) de un buen
orden numerable de tipo \(\beta\). Los códigos se eligen mediante el buen
orden genealógico ya construido. La incrustación ternaria
\[
 E(w)=\sum_{n\geq0}\frac{2\,\mathbf 1_w(n)}{3^{2n+2}}
\]
es inyectiva, y el levantamiento uniforme de cilindros prueba que
\(E(w_\beta)\) pertenece a la recta publicada por HMT. Por tanto,
\[
 k_{\mathrm{HMT}}:
 \omega_1^{\mathfrak G}
 \hookrightarrow\mathbb R_{\mathrm{HMT}}^{\mathfrak G}.
\]
Cantor--Bernstein cierra la igualdad cardinal. La identificación demostrada
\[
 \mathbb R_{\mathrm{HMT}}^{\mathfrak G}
 \approx\mathcal P^{\mathfrak G}(\omega)
\]
la convierte exactamente en \(2^{\aleph_0}=\aleph_1\). Aquí
\(A\approx B\) significa equipotencia; no afirma un isomorfismo de cuerpos
entre una recta y un conjunto potencia.

\clearpage
\subsection{Estructura lógica de la demostración}

La demostración se organiza en cuatro pasos lógicos. Primero se construyen el
dominio de historias, su límite inverso y la cobertura exhaustiva de cilindros.
Segundo se forma la clausura semántica y se aplica la condensación a sus reales.
Tercero se definen, sin usar la igualdad buscada, las dos inyecciones entre la
recta HMT y \(\omega_1^{\mathfrak G}\). Cuarto, Cantor--Bernstein concluye la
igualdad y el argumento de cofinalidad establece la dicotomía para cada
subconjunto interno. En esta secuencia, \(729\) describe la ramificación local;
la cobertura relaciona historias y valores; la condensación controla el nivel
de aparición de los reales internos. Las constantes son publicaciones
correlacionadas del mismo estado: sus cifras convencionales no seleccionan la
comparación cardinal, mientras sus regiones, rutas y memorias permanecen en el
objeto enriquecido sobre el que se construye la recta. La
identificación posterior con \(L[h,m_\infty]\) representa la clausura ya
construida.

El orden expositivo reproduce primero el núcleo formal común de la serie, la
prolongación, la construcción conjunta del continuo y la generación
coinductiva; después, la clausura semántica íntegra y la prueba cardinal. Sólo
entonces desarrolla el núcleo reducido de fase y memoria, el sello
dodecafásico entero \(K\), la vía dodecafásica completa de
\(\alpha_{\mathrm{HMT}}\), el testigo analítico finito de orden nueve, la
representación analítica correlacionada con recurrencia, convergencia y
cilindros comunes, y el corredor excepcional. Estas coordenadas comparten
estado de origen, pero poseen codominios y funciones diferentes; en conjunto
exhiben la doble realización completa y conservan las rutas de publicación, la
incidencia, la frontera y la memoria. Ningún valor objetivo selecciona las
inyecciones cardinales. La comparación con Cantor, Gödel y Cohen se introduce
al término de la construcción~\cite{Cantor1878,Hilbert1902,Godel1938,Cohen1963,Cohen1964}
y no gobierna la definición del objeto HMT.

Cada flecha empleada por el cierre cardinal se define en el cuerpo mediante su
dominio, codominio, acción y compatibilidad con el refinamiento: los tres lazos
orientados del retorno nonádico, su agrupación en bloques de seis símbolos, la
extensión de bloque de \(54\) pasos, la publicación de cilindros y el
levantamiento uniforme de sus ramas. La construcción explícita de las nueve
transiciones proporciona la elevación uniforme empleada en la comparación
cardinal. Los cálculos reproducibles acompañan esta demostración, pero no
sustituyen ninguno de sus argumentos.
